% from ega0-20-fr.tex
\oldpage[0\textsubscript{IV-1}]{132}
Notons en outre que si l'on a un diagramme commutatif d'homomorphismes
d'anneaux
\[
\xymatrix{
  A' \ar[r] & B' \ar[r] & C' \\
  A \ar[u] \ar[r] & B \ar[u] \ar[r] & C \ar[u]
}
\]
on en déduit un diagramme commutatif
\[
\xymatrix{
  \Omega^1_{B/A}\otimes_B C \ar[r] \ar[d] &
  \Omega^1_{C/A} \ar[r] \ar[d] &
  \Omega^1_{C/B} \ar[r] \ar[d] & 0 \\
  \Omega^1_{B'/A'}\otimes_{B'} C' \ar[r] &
  \Omega^1_{C'/A'} \ar[r] &
  \Omega^1_{C'/B'} \ar[r] & 0
}
\tag{20.5.7.3}
\label{0.20.5.7.3-fr}
\]
où les flèches verticales proviennent des di-homomorphismes
\sref{0.20.5.4.3-fr}.

\begin{corollary}[20.5.8]
\label{0.20.5.8-fr}
Supposons que l'homomorphisme $v:B\to C$ fasse de $C$ une $B$-algèbre
formellement étale (pour les topologies discrètes
\sref{0.19.10.2-fr}) ; alors l'homomorphisme \sref{0.20.5.3.3-fr}
\[
  v_{C/B/A}:\Omega^1_{B/A}\otimes_B C\to\Omega^1_{C/A}
\]
est bijectif.
\end{corollary}

\begin{proof}
En effet, si $C$ est une $B$-algèbre formellement non ramifiée pour les
topologies discrètes, il résulte de \sref{0.19.10.4-fr},
\sref{0.20.4.8-fr} et \sref{0.20.1.1-fr} que l'on a
$\operatorname{Hom}_C(\Omega^1_{C/B},L)=0$ pour tout $C$-module $L$, donc
$\Omega^1_{C/B}=0$ (cf. (20.7.4)) ; d'autre part, si $C$ est une
$B$-algèbre formellement lisse pour les topologies discrètes, la suite
\sref{0.20.5.7.2-fr} est exacte ; d'où le corollaire.
\end{proof}

\begin{corollary}[20.5.9]
\label{0.20.5.9-fr}
Soient $A$ un anneau, $B$ une $A$-algèbre, $S$ une partie multiplicative de
$B$ ; alors l'homomorphisme canonique
\begin{equation}
\tag{20.5.9.1}
\label{0.20.5.9.1-fr}
  S^{-1}\Omega^1_{B/A}\to\Omega^1_{S^{-1}B/A}
\end{equation}
est bijectif.
\end{corollary}

\begin{proof}
Il suffit d'appliquer \sref{0.20.5.8-fr} à $C=S^{-1}B$, qui est une
$B$-algèbre formellement étale pour les topologies discrètes
\sref{0.19.10.3-fr}, (ii).
\end{proof}

Compte tenu de \sref{0.20.5.5-fr}, on peut donc écrire
\begin{equation}
\tag{20.5.9.2}
\label{0.20.5.9.2-fr}
  \Omega^1_{S^{-1}B/S^{-1}A}
  =S^{-1}\Omega^1_{B/A}
  =\Omega^1_{S^{-1}B/A},
\end{equation}
à des isomorphismes canoniques près.

\begin{corollary}[20.5.10]
\label{0.20.5.10-fr}
Si $k$ est un corps et $K=k(X_\alpha)_{\alpha\in I}$ une extension pure de
$k$, les $dX_\alpha$ forment une base du $K$-espace vectoriel
$\Omega^1_{K/k}$.
\end{corollary}

\begin{proof}
Comme $K$ est le corps des fractions de l'anneau de polynômes
$k[X_\alpha]_{\alpha\in I}$, cela résulte de \sref{0.20.4.13-fr}, (i) et de
\sref{0.20.5.9-fr}.
\end{proof}

